% !TEX root = main.tex

\section{\texorpdfstring{Lower bound on the complexity of $\beta$-expansions}{}}

\label{sec:lower bound on the complexity of beta expansions}

This section is devoed to the proof of Theorem 1.2. Let $\beta \in (1,2]$. For $s \in [0,1]$, we denote by $\Sigma_\beta(s)$ the set of all $\beta$-expansions of $s$, i.e.,
\begin{equation}
    \Sigma_\beta(s) := \left\{ \xf \in \{0,1\}^\N : \sum_{i=1}^\infty \xf_i \beta^{-i} = s \right\}.
\end{equation}
Recall that, as explain in the first section, $\# \Sigma_2(s) \leq 2$ and $\# \Sigma_\beta(s) = 2^{\aleph_0}$ if $\beta < \frac{1+\sqrt{5}}{2}$, for all $s \in [0,1]$. Further, we define $\Sigma_\beta(s,n)$ to be the set of sequences of length $n$ that are prefix of some $\beta$-expansion of $s$, i.e.,
\begin{equation}
    \Sigma_\beta(s,n) := \{ \xf_{1:n} : \xf \in \Sigma_\beta(s)\}.
\end{equation}

In order to prove Theorem 1.2, we will construct a multivalued function $f_{\beta \to 2} : \{0,1\}^\ast  \rightrightarrows \{0,1\}^\ast$ that is computable relatively to $\beta$, for $\beta \in (1,2]$. This function $f_{\beta \to 2}$ will be designed so that when it is presented with an $n\log_\beta(2)$-prefix $x$ of some $\beta$-expansion of a given real number $s \in [0,1]$, it outputs a set that contains all the $n$-prefixes of all the binary expansions of $s$. In mathematical symbols, this is expressed as
\begin{equation}
    \Sigma_2(s,n) \subseteq f_{\beta \to 2}(x), \forall x \in \Sigma_\beta(s,n\log_\beta(n)), \ n \in \N.
\end{equation}
The function $f_{\beta \to 2}$ will be constructed so that $x \mapsto \# f_{\beta \to 2}(x)$ is bounded, in order to use Theorem \ref{thm:relatively computable multivalued functions and complexity}. 

\subsection{\texorpdfstring{Construction of a multivalued function to convert between bases $\beta$ and $2$}{}}
We now explain the heuristic that is at the origin of the contruction of the function $f_{\beta \to 2}$. This heuristic is inspired from the proof of \cite[Theorem 3]{staiger2002kolmogorov}. Let $\beta \in (1,2]$, $s \in [0,1]$ and $x \in \Sigma_\beta(s,n\log_\beta(2))$. Since $x \in \Sigma_\beta(s,n\log_\beta(2))$, then there exists $\xf \in \{0,1\}^\N$ such that $x\xf$ is a $\beta$-expansion of $s$, i.e.,
\begin{equation}\label{eq:1:construction of a multivalued function computable wrt beta that maps beta expansions to binary expansions of a real number}
    s = \sum_{i=1}^{n\log_\beta(2)} x_i \beta^{-i} + \sum_{i=n\log_\beta(2)+1}^{\infty} \xf_i \beta^{-i}.
\end{equation}
Observe that 
\begin{equation}\label{eq:2:construction of a multivalued function computable wrt beta that maps beta expansions to binary expansions of a real number}
    0 \leq \sum_{i=n\log_\beta(2)+1}^{\infty} \xf_i \beta^{-i} \leq \sum_{i=n\log_\beta(2)+1}^{\infty} \beta^{-i} \leq \frac{\beta^{-n\log_\beta(2)}}{\beta-1} \leq \frac{2^{-n}}{\beta-1}\cdot
\end{equation}
Combining (\ref{eq:1:construction of a multivalued function computable wrt beta that maps beta expansions to binary expansions of a real number}) and (\ref{eq:2:construction of a multivalued function computable wrt beta that maps beta expansions to binary expansions of a real number}), we obtain
\begin{equation}\label{eq:3:construction of a multivalued function computable wrt beta that maps beta expansions to binary expansions of a real number}
    \sum_{i=1}^{n\log_\beta(2)} x_i \beta^{-i} \leq s \leq \sum_{i=1}^{n\log_\beta(2)} x_i \beta^{-i} + \frac{2^{-n}}{\beta-1}\cdot
\end{equation}
Therefore, from the knowledge of $x$ only, we can deduce that $s$ belongs to the interval $I(x)$ defined as 
\begin{equation}\label{eq:4:construction of a multivalued function computable wrt beta that maps beta expansions to binary expansions of a real number}
    I(x) :=  \left[\sum_{i=1}^{n\log_\beta(2)} x_i \beta^{-i},  \sum_{i=1}^{n\log_\beta(2)} x_i \beta^{-i}+ \frac{2^{-n}}{\beta-1} \right].
\end{equation}
Now, by combining again (\ref{eq:1:construction of a multivalued function computable wrt beta that maps beta expansions to binary expansions of a real number}) and (\ref{eq:2:construction of a multivalued function computable wrt beta that maps beta expansions to binary expansions of a real number}) for the specific case of $\beta = 2$, we get that
\begin{equation}\label{eq:5:construction of a multivalued function computable wrt beta that maps beta expansions to binary expansions of a real number}
    t - 2^{-n} \leq \sum_{i=1}^{n} y_i 2^{-i} \leq t,
\end{equation}
for all $t \in [0,1]$ and $y \in \Sigma_2(t,n)$. This means that if we know exactly the value of $t$, we can deduce that all the $n$-prefixes of the binary expansions of $t$ are exactly those sequences $y \in \{0,1\}^n$ that satisfy 
\begin{equation}\label{eq:6:construction of a multivalued function computable wrt beta that maps beta expansions to binary expansions of a real number}
    \sum_{i=1}^{n} y_i 2^{-i} \in [t-2^{-n}, t].
\end{equation}
However, if we do not know exactly $t$, but rather we know that $t \in [a,b] \subseteq [0,1]$, then we can deduce that all the $n$-prefixes $y$ of the binary expansions of $t$ satisfy
\begin{equation}\label{eq:7:construction of a multivalued function computable wrt beta that maps beta expansions to binary expansions of a real number}
    \sum_{i=1}^{n} y_i 2^{-i} \in [a-2^{-n}, b].
\end{equation}
Now, in light of (\ref{eq:7:construction of a multivalued function computable wrt beta that maps beta expansions to binary expansions of a real number}), (\ref{eq:4:construction of a multivalued function computable wrt beta that maps beta expansions to binary expansions of a real number}) can be reinterpreted as follows. From the knowledge of the $x$, we can deduce that all the $n$-prefixes $y$ of the binary expansions of $s$ satisfy
\begin{equation}\label{eq:J x:construction of a multivalued function computable wrt beta that maps beta expansions to binary expansions of a real number}
    \sum_{i=1}^{n} y_i 2^{-i} \in \left[\sum_{i=1}^{n\log_\beta(2)} x_i \beta^{-i} - 2^{-n},  \sum_{i=1}^{n\log_\beta(2)} x_i \beta^{-i}+ \frac{2^{-n}}{\beta-1} \right] =: J(x).
\end{equation}
The above heuristic results in the following formal definition of a multivalued function $\hat f_{\beta \to 2}$, that however is still not the multivalued function $f_{\beta \to 2}$ we are looking for, for reasons we explain just after the staement of the definition.
\begin{definition}\label{def:noncomputable multivalued function converting from base beta to base 2}
    Let $\hat f_{\beta \to 2} : \{0,1\}^\ast \rightrightarrows \{0,1\}^\ast$ be the multivalued function defined as 
\begin{equation}\label{eq:main:def:noncomputable multivalued function converting from base beta to base 2}
    \hat f_{\beta \to 2}(x) := \left\{ y \in \{0,1\}^n :  \sum_{i=1}^{n} y_i 2^{-i} \in J(x)\right\},
\end{equation}
for all $x \in \{0,1\}^{n\log_\beta(2)}$ for some $n \in \N$, and $\hat f_{\beta \to 2}(x) = \epsilon$ for $x \in \{0,1\}^m$ where $m \neq n \log_\beta(2)$ for all $n \in \N$.
\end{definition} 
However, this function is clearly not computable (even relatively) since it makes use of comparisons between arbitrary real numbers to evalute wether or not $\sum_{i=1}^{n} y_i 2^{-i} \in J(x)$, i.e., its output cannot be calculated from its input through an effective algorithm. The core of the problem is that the operation ``$\in$'' is not computable, as it relies on both the operations ``$\leq$'' and ``$\geq$''. This well-known issue can be fixed by using an approximate version of ``$\leq$''. Namely, following \cite[Section 4.1, p79-80]{brattka2018ComputabilityAnalysis}, there exists a computable function $\varphi : \N \times \{0,1\}^\N \to \{0,1\}$ such that
\begin{equation}
    \varphi(n, \langle \yf, \zf \rangle) = \begin{cases}
        0 \ \text{if} \ s < t\\
        0 \ \text{or} \ 1 \ \text{if} \ t \leq s \leq t + 2^{-n},\\
        1 \ \text{if} \ t > s + 2^{-n},
    \end{cases}
\end{equation}
for all $n \in \N$, $s,t \in [0,1]$, and where $\yf$ and $\zf$ are the respective greedy binary expansions of $s$ and $t$. For $n \in \N$, we define the binary relation $\leq_n$ by 
\begin{equation}
    s \leq_n t \iff \varphi(n , \langle \yf, \zf \rangle), \ \forall s,t \in [0,1],
\end{equation}
where $\yf,\zf$ are the respective greedy binary expansions of $s$ and $t$. Based on this relation, we can construct an ``approximate membership'' relation for $n \in \N$ by 
\begin{equation}
    s \in_n [a,b] \iff a \leq_n s - 2^{-n} \ \text{and} \ a\leq_n b, \ \forall s,a,b \in [0,1].
\end{equation}
Note that, in particular,
\begin{equation}
    s \in_n [a,b] \Rightarrow s \in [a - 2^{-n}, b + 2^{-n}], \ \ \forall s,a,b \in [0,1].
\end{equation}
In consequence, for an interval $I := [a,b] \subseteq [0,1]$, we define 
\begin{equation}\label{eq:definition of superset of size 2 -n more}
    I^{(n)} := [a - 2^{-n}, b + 2^{-n}], \ \text{for} \ n \in \N,
\end{equation}
and we get 
\begin{equation}\label{eq:finite precision membership up to 2 -n implies membership to a superset of size 2 -n more}
    s \in_n I \Rightarrow s \in I^{(n)}, \ \ \forall s \in [0,1].
\end{equation}
Now, the function $f_{\beta \to 2}$ can simply be defined with the usage of this approxiamte membership relation.
\begin{definition}\label{def:multivalued function converting from base beta to base 2}
    Let $f_{\beta \to 2} : \{0,1\}^\ast \rightrightarrows \{0,1\}^\ast$ be the multivalued function defined as 
\begin{equation}\label{eq:main:def:multivalued function converting from base beta to base 2}
    f_{\beta \to 2}(x) := \left\{ y \in \{0,1\}^n :  \sum_{i=1}^{n} y_i 2^{-i} \in_n J(x)\right\},
\end{equation}
for all $x \in \{0,1\}^{n\log_\beta(2)}$ for some $n \in \N$, and $f_{\beta \to 2}(x) = \epsilon$ for $x \in \{0,1\}^m$ where $m \neq n \log_\beta(2)$ for all $n \in \N$.
\end{definition} 

We construct another multivalued function $f_{\beta \to 2}$ as a slight modification of $\hat f_{\beta \to 2}$, which however is computable with respect to $\beta$. To do so, we first need to develop a systematic way to approximate comparison between real numbers with computable functions.


%% ATTENTION ATTENTION ATTENTION
% NE PAS ENLEVER, CETTE SECTIOON CONTIENT DES TRUCS STYLES
% \subsection{Computable approximate threshold functions}

% \begin{definition}
%     For all $s \in \R$, and $n \in \N$ we define 
%     \begin{equation}
%         s^{(n)} := \lfloor s \rfloor + \sum_{i=1}^n \yf_i 2^{-i},
%     \end{equation}
%     where $\yf \in \{0,1\}^\N$ is the greedy binary expansion of $\{s\} := s - \lfloor s \rfloor \in [0,1)$.
% \end{definition}
% Note that $s - s^{(n)} \leq 2^{-n}$, for all $s \in \R$, $n \in \N$. We define an approximate ordering relationship between real numbers.
% \begin{definition}
%     For $m \in \N$, we define the binary relation $\leq_m$ over $\R \times \R$ as 
%     \begin{equation}
%         s \leq_m t \iff s^{(m)} \leq t^{(m)} + 2^{-m},
%     \end{equation}
%     for all $s,t \in \R$.
% \end{definition}
% We recall here the definition of \textit{rational functions}. A rational function over $\R$ is a function of the form
% \begin{equation}
%     \begin{array}{lrcl}
%         g : & \dom g & \to & \R\\
%             & x & \mapsto & \frac{p_1(x)}{p_2(x)},
%     \end{array}
% \end{equation}
% where $p_1,p_2 : \R \to \R$ are two polynomials, and $\dom g := \R \backslash \{ x \in \R: p_2(x) = 0\}$. Note that then, if $p_2$ has $k$ roots, we denote by $r_1 < \ldots < r_k$ its roots, and then $\dom g = \R \backslash \{r_1, \ldots, r_k\}$. With $r_0 := -\infty$ and $r_{k+1} := \infty$, we then have 
% \begin{equation}
%     \dom g = \bigcup_{i=1}^k (r_i,r_{i+1}).
% \end{equation}
% An open interval $I = (a,b) \subseteq \R$ is called an \textit{interval of continuity} of $g$ if $g$ is continuous on $I$. Note that, for $i \in \{0,\ldots, k\}$, $(r_i, r_{i+1})$ is an interval of continuity of $g$, since $g$ is obviously continuous on such an interval. 
% % For a family $\{g_a\}_{a \in A}$ of rational functions, where $A$ is some set, we say that an open interval $I= (a,b) \subseteq \R$ is a common interval of continuity of $\{g_a\}_{a \in A}$ if for every $a \in A$, $g_a$ has an interval of continuity $J$ such that $I \subseteq J$. 
% We consider the class $\Q(X)$ of rational functions for which that the coefficients of the polynomials $p_1$ and $p_2$ are in $\Q$. 
% \begin{lemma}
%     Let $f \in \Q(X)$. Then, there exists a computable function $\psi : \{0,1\}^\ast \to \Q$ such that 
%     \begin{equation}
%         \psi_f \left(x\right) = f\left(\sum_{i=1}^{|x|} x_i 2^{-i}\right),
%     \end{equation}
%     for all $x \in \{0,1\}^\ast$.
% \end{lemma}
% \begin{proof}
%     Let $f \in \Q(X)$. Then, there exists $k,\ell \in \N$ and $a_0,\ldots, a_k, b_0, \ldots, b_\ell \in \Q$ such that 
%     \begin{equation}
%         f(s) = \frac{\sum_{j=0}^k a_j s^j}{\sum_{j=0}^\ell b_j s^j}, \ \forall s \in \R.
%     \end{equation}
%     Then, for all $x \in \{0,1\}$,
%     \begin{equation}
%         f\left(\sum_{i=1}^{|x|} x_i 2^{-i}\right) = \frac{\sum_{j=0}^k a_j \left(\sum_{i=1}^{|x|} x_i 2^{-i}\right)^j}{\sum_{j=0}^\ell b_j \left(\sum_{i=1}^{|x|} x_i 2^{-i}\right)^j}\cdot
%     \end{equation}
%     Since ABOVE EQUATION only consists of additions, multiplications and divisions between rational numbers, the exists an effective algorithm that convert any $x \in \{0,1\}^\ast$ to $f\left(\sum_{i=1}^{|x|} x_i 2^{-i}\right)$. This effective algorithm is underlying a computable function $\psi_f$.
% \end{proof}
% \begin{lemma}
%     Let $f \in \Q(X)$ and suppose that $(0,1)$ is an interval of continuity of $f$. Then, there exists a computable function $\varphi_f : \N \to \Q$ such that 
%     \begin{equation}
%         \left|\varphi_f^{\yf} \left(n\right) - f(s)\right| \leq 2^{-n},
%     \end{equation}
%     for all $s \in (0,1)$, where $\yf$ is the greedy binary expansion of $s$.
% \end{lemma}
% \begin{proof}
%     Let $f \in \Q(X)$ and suppose that $(0,1)$ is an interval of continuity of $f$. Define $p_1,p_2 \in \Q[X]$ such that $f(s) = p_1(s)/p_2(s)$, for all $s \in (0,1)$. We will make use of the last LEMMA in combination with continuity arguments. Namely, for a given $s \in (0,1)$, its greedy binary expansion $\yf \in \{0,1\}^\N$ can be used to approximate $s$ arbitrarily with rational numbers. Namely, for $m \in \N$, define $s^{(m)} := \sum_{i=1}^m \yf_i 2^{-i}$. By previous LEMMA, there exists a computable function $\psi_f: \{0,1\}^\ast \to \Q$ such that $\psi_f \left(x\right) = f\left(\sum_{i=1}^{|x|} x_i 2^{-i}\right)$. In particular, $\psi(\yf_{1:m}) = f(s^{m})$. By continuity of $f$, for a given $n \in \N$, there exists $m_n \in \N$, such that $|f(s) - f(s^{(m)})| \leq 2^{-n}$ for all $m \geq m_n$. Suppose that there exists a computable function $\eta_f : \N \to \N$ such that $\eta_f(n) \geq m_n$ for all $n \in \N$. Then, we can define $\varphi_f : \N \to \Q$ such that
%     \begin{equation}
%         \varphi_f^\yf(n) := \psi_f(\yf_{1:\eta_f(n)}), \ \forall n \in \N.
%     \end{equation}
%     We now turn to the construction of $\eta_f : \N \to \N$.
%     \begin{enumerate}[label = (\alph*)]
%         \item Suppose that $f$ is a polynomial, i.e., $f \in \Q[X] \subseteq \Q(X)$. Let $d$ be the degree of $f$, and $a_0, \ldots, a_d$ be the coeffcients of $f$, such that $f(s) = \sum_{i=0}^d a_i s^i$, for all $s \in (0,1)$. Note that
%         \begin{align}
%             |f(s) - f(s^{(m)})| &\leq \sum_{i=0}^d |a_i| \left(s^i - \left(s^{(m)}\right)^i\right)\\
%                 &\leq  \sum_{i=0}^d |a_i| \left((s^{(m)}+2^{-m})^i - \left(s^{(m)}\right)^i\right)\\
%                 &= \sum_{i=0}^d |a_i| \sum_{j=1}^i \binom{i}{j} (s^{(m)})^j 2^{-m(i-j)}\\
%                 &\leq \sum_{i=0}^d |a_i| \sum_{j=1}^i \binom{i}{j}  2^{-m(i-j)}\\
%                 &= \sum_{i=0}^d |a_i| \left((1+2^{-m})^i - 1\right) =: \varepsilon_m,
%         \end{align}
%         for all $s \in (0,1)$, where PUT EXPLANATIONS. Note that $\limi{m} \varepsilon = 0$. Hence, we construct an effective algorithm for $\eta_f : \N \to \N$ as follows. On input $n \in \N$, calculate $\varepsilon_1, \varepsilon_2, \ldots$ until reaching $m \in \N$ such that $\varepsilon_m \leq 2^{-n}$, and output $m \in \N$. This concludes the construction of $\eta_f : \N \to \N$, for $f$ polynomial. 
%         \item We now move to the general case $f \in \Q(X)$. Recall that since $f \in \Q(X)$, there exists two polynomials $p_1,p_2 \in \Q[X]$ such that $f(s) = \frac{p_1(s)}{p_2(s)}$, for all $s \in (0,1)$, and $p_2(s) \neq 0$ for all $s \in (0,1)$ (since by assumption, $(0,1)$ is an interval of continuity of $f$). Note that 
%         \begin{align}
%             |f(s) - f(s^{(m)})| &\leq (s - s^{(m)}) \max_{t \in [s^{(m)}, s]} |f'(s)|\\
%                 &\leq 2^{-m} \max_{t \in [s^{(m)}, s]} \frac{|p'_1(t)p_2(t) - p_1(t)p'_2(t)|}{(p_2(t))^2}\\
%                 &\leq 2^{-m}  \frac{\displaystyle \max_{t \in [s^{(m)}, s]} \left|p'_1(t)p_2(t) - p_1(t)p'_2(t)\right|}{\displaystyle \min_{t \in [s^{(m)}, s]}(p_2(t))^2}\cdot
%         \end{align}
%         For $m \in \N$, $q_1 := p'_1p_2 - p_1p'_2$ and $q_2 := p_2^2$ are both polynomials with rational coefficients. Then, we can use the computable functions $\varphi_{q_1}, \varphi_{q_2} : \N \to \Q$ to compute $M_m, N_m > 0$ such that 
%         \begin{equation}
%             \frac{M_n}{2} \leq \max_{t \in [s, s^{(m)}]} \left|q_1(t)\right| \leq M_m \ \ \text{and} \ \ N_m \leq \min_{t \in [s, s^{(m)}]} q_2(t) \leq 2 N_m.
%         \end{equation}
%         Therefore, we have 
%         \begin{equation}
%             |f(s) - f(s^{(m)})| \leq 2^{-m} \frac{M_n}{N_m} =: \varepsilon_m.
%         \end{equation}
%         Note that $\limi{m} \varepsilon = 0$. Hence, we construct an effective algorithm for $\eta_f : \N \to \N$ as follows. On input $n \in \N$, calculate $\varepsilon_1, \varepsilon_2, \ldots$ until reaching $m \in \N$ such that $\varepsilon_m \leq 2^{-n}$, and output $m \in \N$. This concludes the proof. 
%     \end{enumerate}
% \end{proof}
% We can now use the concept of rational functions and the approximate ordering relationship to construct a relatively computable function. 
% \begin{theorem}
%     Consider a family of functions $(f_x)_{x \in \{0,1\}^\ast} \in \Q(X)$, $g \in \Q(X)$, and suppose that $(0,1)$ is an interval of continuity for $g$ and $f_x$, for all $x \in \{0,1\}^\ast$. Define the function $h : (0,1) \times \{0,1\}^\ast \times \N \to \{0,1\}$ by
%     \begin{equation}
%         h(s, x, m ) := \begin{cases}
%             0 \ if \ f_x(s) \leq_m g(s),\\
%             1 \ if \ f_x(s) >_m g(s).
%         \end{cases}
%     \end{equation}
%     Then, there exists a computable function $\varphi_h : \{0,1\}^\ast \times \N \to \{0,1\}$ such that 
%     \begin{equation}
%         \chi^\yf(x ,m) = h(s,x,m), \ \forall s \in (0,1), \ x \in \{0,1\}^\ast, \ \text{and} \ m \in \N,
%     \end{equation}
%     where $\yf \in \{0,1\}^\N$ is the greedy binary expansion of $s$.
% \end{theorem}
% \begin{proof}
%     Consider a family of functions $(f_x)_{x \in \{0,1\}^\ast} \in \Q(X)$, $g \in \Q(X)$, and suppose that $(0,1)$ is an interval of continuity for $g$ and $f_x$, for all $x \in \{0,1\}^\ast$. Let $s \in (0,1)$. We construct an effective algorithm for $\chi$ as follows. Let $s \in (0,1)$, $x \in \{0,1\}^\ast$, $m \in \N$, and suppose that we equip the algorithm with the greedy binary expansion $\yf$ of $s$. Then, when presented with input $(x,m)$, the algorithm computes $a_m := \varphi_{f_x}(m+1)$ and $b_m := \varphi_g(m+1)$, where $\varphi_{f_x},\varphi_{g}$ are the functions defined by LEMMA. Then, the algorithm outputs 
% \end{proof}
% We further define an approximate membership relation based on the approximate relation we just derived.
% \begin{definition}
%     For $m \in \N$, we define the binary relation $\in_m$ over $\R^3$ as 
%     \begin{equation}
%         s \in_m [a,b] \iff a \leq_m s \leq_m b,
%     \end{equation}
%     for all $s,a,b \in \R$.
% \end{definition} 
% \begin{corollary}
%     Consider a family of continuous functions $(f_x : \R \to \R)_{x \in \{0,1\}^\ast}$, and two continuous functions $g_1,g_2 : \R \to \R$. Then, for $s \in \R$, the function $h_s : \{0,1\}^\ast \to \{0,1\}^\ast$ defined by 
%     \begin{equation}
%         h_s(\langle x,m \rangle) := \begin{cases}
%             0 \ if \ f_x(s) \in_m [g_1(s), g_2(s)]\\
%             1 \ if \ f_x(s) \notin_m [g_1(s), g_2(s)]
%         \end{cases}
%     \end{equation}
%     is computable relatively to $s$, for all $s \in \R$.
% \end{corollary}
% \begin{proof}
%     This is a direct consequence of THEOREM.
% \end{proof}
% As a consequence, we define a modification of $e_{\beta \to 2}$ as follows.
% \begin{definition}\label{def:computable multivalued function converting from base beta to base 2}
%     Let $f_{\beta \to 2} : \{0,1\}^\ast \rightrightarrows \{0,1\}^\ast$ be the multivalued function defined as 
% \begin{equation}\label{eq:main:def:computable multivalued function converting from base beta to base 2}
%     f_{\beta \to 2}(x) := \left\{ y \in \{0,1\}^n :  \sum_{i=1}^{n} y_i 2^{-i} \in_n J(x)\right\},
% \end{equation}
% for all $x \in \{0,1\}^{n\log_\beta(2)}$ for some $n \in \N$, and $f_{\beta \to 2}(x) = \epsilon$ for $x \in \{0,1\}^m$ where $m \neq n \log_\beta(2)$ for all $n \in \N$.
% \end{definition} 

\subsection{Consequences on algorithmic complexity}


 This function is clearly computable relatively to $\beta$, since we can design an algorithm that queries $x$ and bits of the greedy binary expansion of $\beta-1$, until it can enumerate all the $y \in \{0,1\}^n$ that satisfy (\ref{eq:J x:construction of a multivalued function computable wrt beta that maps beta expansions to binary expansions of a real number}). We now establish an upper bound on the cardinality of $f_{\beta \to 2}(x)$, for all $x \in \{0,1\}^\ast$.
\begin{lemma}\label{lem:upper bound on the cardinality of the beta to 2 conversion multifunction}
    For $\beta \in (1,2]$ and $x \in \{0,1\}^\ast$,
    \begin{equation}
        \#f_{\beta \to 2}(x) \leq \frac{1}{\beta-1} + 3, \ \forall x \in \{0,1\}^\ast.
    \end{equation}
\end{lemma}
\begin{proof}
    Let $\beta \in (1,2]$ and $x \in \{0,1\}^\ast$. 
    \begin{enumerate}[label=(\alph*)]
        \item Suppose that there is no $n \in \N$ such that $\len x = n\log_\beta(2)$. Then, $f(x) = \epsilon$, so $\#f(x) = 1 \leq \frac{1}{\beta-1} + 3$.
        \item Suppose that there exists $n \in \N$ such that $\len x = n\log_\beta(2)$. By definition,
        \begin{align}
            f_{\beta \to 2}(x) &= \left\{ y \in \{0,1\}^n :  \sum_{i=1}^{n} y_i 2^{-i} \in_n J(x)\right\}\\
            &\subseteq \left\{ y \in \{0,1\}^n :  \sum_{i=1}^{n} y_i 2^{-i} \in J(x)^{(n)}\right\} =: A(x)
        \end{align}
        We give an upper bound on the cardinality of $A(x)$, which turns out to be an upper bound on the cardinality of $f_{\beta \to 2}(x)$. Note that for $y,z \in \{0,1\}^n$, satisfying $y \neq z$, we have
        \begin{equation}
            \left|\sum_{i=1}^{n} y_i 2^{-i} - \sum_{i=1}^{n} z_i 2^{-i}\right| \geq 2^{-n}.
        \end{equation}
        Therefore, for an interval $I = [a,b]$, there can be at most $2^n(b-a)$ sequences $y \in \{0,1\}^n$ that satisfy $\sum_{i=1}^{n} y_i 2^{-i} \in I$. 
        By (\ref{eq:J x:construction of a multivalued function computable wrt beta that maps beta expansions to binary expansions of a real number}) and (\ref{eq:definition of superset of size 2 -n more}),
        \begin{equation}
            J(x)^{(n)} = \left[\sum_{i=1}^{n\log_\beta(2)} x_i \beta^{-i} - 2^{-n} - 2^{-n}, \sum_{i=1}^{n\log_\beta(2)} x_i \beta^{-i}+ \frac{2^{-n}}{\beta-1} + 2^{-n}\right],
        \end{equation}
        so
        \begin{align}
            \#A(x) &\leq 2^n \left(\sum_{i=1}^{n\log_\beta(2)} x_i \beta^{-i}+ \frac{2^{-n}}{\beta-1} + 2^{-n} - \left(\sum_{i=1}^{n\log_\beta(2)} x_i \beta^{-i} - 2^{-n} - 2^{-n} \right)\right)\\
            &= 2^n \left(2^{-n}\frac{1}{\beta-1} + 3\cdot2^{-n}\right) = \frac{1}{\beta-1} + 3.
        \end{align}
        As $f_{\beta \to 2}(x) \subseteq A(x)$, we finally get
        \begin{equation}
            \#f_{\beta \to 2}(x) \leq \#A(x) \leq \frac{1}{\beta-1} + 3.
        \end{equation}
    \end{enumerate}
\end{proof}
\noindent Finally, we prove Theorem \ref{thm:binary expansions are at least as compressible as as beta expansions} as a corollary of Theorem \ref{thm:relatively computable multivalued functions and complexity} and Lemma \ref{lem:upper bound on the cardinality of the beta to 2 conversion multifunction}. 

\begin{corollary}\label{cor:binary expansions are at least as compressible as as beta expansions}
    Let $\beta \in (1,2)$. Then,
		\begin{equation}
			0 \leq \underline\Delta_\beta(\xf) \leq \bar\Delta_\beta(\xf),
		\end{equation}
		for all $\xf \in \Omega_\beta$.
\end{corollary}
\begin{proof}
    Let $\beta \in (1,2)$, $\xf \in \Omega_\beta$, and $s := \sum_{i=1}^\infty \xf_i \beta^{-i}$. Then, by Lemma \ref{lem:upper bound on the cardinality of the beta to 2 conversion multifunction}, there exists a multivalued function $f : \{0,1\}^\ast  \rightrightarrows \{0,1\}^\ast$ that is computable relatively to $\beta$, such that
    \begin{equation}
        \Sigma_2(s,n) \subseteq f(\xf_{1:n\log_\beta(2)}),
    \end{equation}
    and
    \begin{equation}
        \#f(\xf_{1:n\log_\beta(2)}) \leq \frac{1}{\beta-1} + 3.
    \end{equation}
    Consequently, Theorem \ref{thm:relatively computable multivalued functions and complexity} yields that 
    \begin{align}
        K[y] &\leq K[\xf_{1:n\log_\beta(2)}] + \log \left(\frac{1}{\beta-1} + 3\right) + \log\log \left(\frac{1}{\beta-1} + 3\right) + \bigOx[1]{n}\\
        &= K[\xf_{1:n\log_\beta(2)}] + \bigOx[1]{n},
    \end{align}
    for all $y \in  \Sigma_2(s,n)$. In particular, if $\yf \in \{0,1\}^\N$ is the greedy binary expansion of $s$, then 
    \begin{equation}
        K[\yf_{1:n}] \leq K[\xf_{1:n\log_\beta(2)}] + \bigOx[1]{n}.
    \end{equation}
    Finally, this yields
    \begin{equation}
        \bar{\Delta}_{\beta}(\xf) \geq \underline{\Delta}_{\beta}(\xf) = \liminf_{n\to \infty} \frac{K[\xf_{1:n\log_\beta(2)}] - K[\yf_{1:n}]}{n} \geq \liminf_{n\to \infty}\bigOx[n^{-1}]{n} = 0.
    \end{equation}
\end{proof}
